\chapter{Antecedentes}

A continuación, se presentan los principales antecedentes relacionados con el estudio de esta propuesta doctoral, enfocados en la problemática de gestión de residuos, el potencial de soluciones biotecnológicas y el papel de tecnologías emergentes para mitigar el impacto ambiental.

\section{Problemática global en el manejo de residuos y la generación de emisiones de GEI}

La gestión inadecuada de residuos constituye una de las principales amenazas ambientales a escala global, con efectos directos sobre el cambio climático, la salud pública y la integridad de los ecosistemas. Según la FAO, se estima que más del 50\,\% de los residuos sólidos generados en países de ingresos medios y bajos son de origen orgánico \parencite{FAO2021}, los cuales, al ser dispuestos sin tratamiento adecuado, se degradan en condiciones anaeróbicas generando metano (CH$_4$), un gas con un potencial de calentamiento global 28 veces mayor que el CO$_2$ a 100 años. Este fenómeno es particularmente grave en sectores como la agricultura, donde las prácticas tradicionales —como la quema de residuos y el uso intensivo de insumos químicos— agravan la carga ambiental de los sistemas productivos.

De acuerdo con el Quinto Informe del Panel Intergubernamental sobre Cambio Climático, aproximadamente el 24\% de las emisiones globales de GEI provienen de actividades relacionadas con el uso del suelo, la agricultura y los cambios en el uso de la tierra \parencite{IPCC2014}. A estas emisiones se suman los impactos generados por el manejo ineficiente de residuos en contextos urbanos e industriales, incluyendo lixiviación contaminante, contaminación atmosférica y afectación de la biodiversidad local. Estas condiciones se ven agravadas por la limitada infraestructura de tratamiento y las brechas en gobernanza ambiental, especialmente en regiones rurales y en países en desarrollo.

Este contexto exige de manera urgente un diseño de estrategias sostenibles de gestión de residuos, que integren enfoques sistémicos y soluciones basadas en principios de economía circular. El enfoque de “residuo como recurso” promueve la valorización de la materia orgánica a través de tecnologías como el compostaje, la digestión anaerobia, la vermicultura y, más recientemente, la bioconversión con insectos. Estos enfoques permiten cerrar ciclos de nutrientes, reducir emisiones y generar productos útiles como biofertilizantes, energía o proteínas alternativas, entre otros.

No obstante, la aplicación masiva de estas soluciones aún enfrenta barreras tecnológicas, económicas y de medición. Para lograr una transformación efectiva, es necesario avanzar hacia sistemas de gestión inteligentes, capaces de monitorear en tiempo real la generación de emisiones y otras variables críticas. Esto exige la incorporación de tecnologías emergentes como sensores con Internet de las Cosas (IoT), modelos matemáticos de predicción y algoritmos de aprendizaje automático, que permitan mejorar la eficiencia de los sistemas, garantizar la trazabilidad de los procesos y sustentar decisiones basadas en datos confiables \parencite{ahmedAgricultureClimateChange2020}.

\section{Estudios Previos sobre la Gestión de Residuos con \textit{BSF}}

En los últimos años, numerosos estudios han destacado el potencial de \textit{BSF} como una solución biotecnológica eficaz para la gestión de residuos orgánicos. Su alta eficiencia en la conversión de materia orgánica en biomasa larvaria rica en proteínas y lípidos, así como en fertilizante orgánico (frass), ha sido ampliamente documentada. Por ejemplo, \textcite{salamEffectDifferentEnvironmental2022} evaluaron el efecto de variables ambientales sobre la eficiencia del proceso, demostrando que temperaturas entre $25-30\ ^{\circ}\mathrm{C}$ y niveles de humedad superiores al $60\%$ son condiciones óptimas para maximizar la tasa de bioconversión. Estas observaciones han sido clave para definir parámetros operativos en sistemas de cría intensiva.

Asimismo, estudios recientes han ampliado el espectro de aplicaciones de la biomasa generada. Investigaciones como las de \textcite{koyunogluBiofuelProductionUtilizing2024} y \textcite{suryatiLauricAcidBlack2023} han explorado su uso en la producción de biocombustibles y bioplásticos, así como en la obtención de ácidos grasos de cadena media para alimentación animal o incluso humana. Estos avances han reforzado el posicionamiento de \textit{BSF} dentro de los modelos de economía circular, al permitir la valorización integral de residuos y la generación de productos con alto valor agregado.

Sin embargo, la mayoría de estas investigaciones se han enfocado en aspectos aislados del proceso -como la optimización de condiciones ambientales, la cinética de crecimiento larval o la caracterización de subproductos- sin abordar de forma holística su impacto ambiental ni su potencial para ser integrados en sistemas inteligentes de gestión. En particular, son escasos los estudios que incorporan herramientas de modelado matemático para predecir la generación de gases de efecto invernadero (GEI) durante la bioconversión, o que integren tecnologías emergentes como sensores con IoT y algoritmos de aprendizaje automático (AAA) en la evaluación operativa del sistema.

\section{Medición de Emisiones de Gases de Efecto Invernadero en la Bioconversión con \textit{BSF}}

La medición de gases de efecto invernadero (GEI) es un componente esencial para evaluar la sostenibilidad ambiental de los sistemas de bioconversión basados en \textit{BSF}. Estos procesos, aunque eficientes en la transformación de residuos orgánicos en biomasa útil, generan emisiones significativas de CO$_2$, CH$_4$ y N$_2$O debido a la actividad metabólica de las larvas y la descomposición microbiana del sustrato. Estudios recientes han empleado técnicas como cromatografía de gases, sensores NDIR (Non-Dispersive Infrared) y sensores electroquímicos para cuantificar estas emisiones, tanto en condiciones experimentales como operativas \parencite{galan-diazCarbonWaterFootprint2024,rossiEstimatingDynamicsGreenhouse2024,bekkerImpactSubstrateMoisture2021a}.

Por ejemplo, \textcite{jenkinsProcessingPoultryManure2025} demostraron que los sistemas de bioconversión con \textit{BSF} pueden reducir significativamente las emisiones de CO$_2$ y N$_2$O  en comparación con métodos convencionales como el compostaje, debido a un menor tiempo de permanencia del residuo y una actividad microbiana más controlada. Asimismo, \textcite{fuhrmannComprehensiveIndustryrelevantBlack2025} reportaron valores promedio de $96.5\ \text{g CO}_2$ por kg de residuo tratado, mientras que \textcite{rossiEstimatingDynamicsGreenhouse2024} observaron entre $7.76$ y $11.88\ \text{g CO}_2$ por gramo de larva seca, subrayando la necesidad de contextualizar las emisiones según el tipo de medición y unidad funcional empleada.

No obstante, la precisión y comparabilidad de estas mediciones se ven afectadas por varios factores críticos, como la calibración de los sensores, las condiciones ambientales durante el muestreo, la composición del sustrato y la representatividad de las unidades experimentales. La heterogeneidad en los protocolos de medición, la escasa estandarización metodológica y la falta de reportes de incertidumbre limitan la capacidad de comparar resultados entre estudios y, por ende, dificultan el diseño de estrategias robustas de mitigación.

Además, la literatura presenta una notoria carencia de estudios en condiciones a gran escala o semicomerciales, lo que restringe la extrapolación de resultados hacia escenarios industriales o regionales. Esta brecha evidencia la necesidad de desarrollar plataformas tecnológicas de medición continua, integradas con sensores con IoT y sistemas de análisis de datos en tiempo real, que permitan caracterizar de forma dinámica y precisa las emisiones de GEI en distintos entornos de bioconversión. La incorporación de sistemas de Medición, Reporte y Verificación (MRV) basados en tecnologías digitales puede representar un salto cualitativo en la validación ambiental de estos procesos \parencite{kamilaris2017review,zhangApplicationBigData2024,rajRealTimeEstimation2023}.

\section{Antecedentes en Modelamiento Matemático de Procesos de Bioconversión}

El modelamiento matemático ha sido una herramienta clave para comprender, simular y optimizar procesos biológicos complejos relacionados con la gestión de residuos. En el contexto de la bioconversión mediada por insectos, como \textit{BSF}, los modelos matemáticos permiten describir de manera cuantitativa la transformación de la materia orgánica en biomasa y subproductos, así como las emisiones asociadas de gases de efecto invernadero (GEI), como el dióxido de carbono (CO$_2$) y el metano (CH$_4$).

Entre los enfoques más utilizados se encuentran los modelos cinéticos no lineales y los sistemas de ecuaciones diferenciales ordinarias (EDO), que describen la dinámica de variables como el crecimiento larval, la degradación del sustrato y la generación de gases. Por ejemplo, \textcite{padmanabhaComprehensiveDynamicGrowth2020} desarrollaron un modelo basado en el marco del Dynamic Energy Budget (DEB), dividiendo la biomasa larval en compartimentos estructurales, lipídicos y de asimilados, lo que permitió simular los flujos metabólicos y las emisiones respiratorias. De manera complementaria, \textcite{eriksenDynamicModellingFeed2022} propuso un modelo compartimental que incorpora balances energéticos y principios logísticos para describir la acumulación de lípidos y la producción de CO$_2$ en condiciones de cría controlada.

Además, otros estudios han utilizado funciones tipo Monod y Haldane para representar las tasas de consumo de sustrato y la actividad microbiana en contextos de digestión anaerobia o vermicompostaje, los cuales comparten similitudes funcionales con los procesos de bioconversión con \textit{BSF} \parencite{winAnaerobicDigestionBlack2018, ahlamineMathematicalAnalysisAnaerobic2024}. Estas aproximaciones permiten capturar las no linealidades inherentes al sistema y evaluar cómo variables ambientales como la temperatura, la humedad y la disponibilidad de oxígeno modulan la generación de GEI \parencite{bekkerImpactSubstrateMoisture2021a, rossiEstimatingDynamicsGreenhouse2024}.

No obstante, la literatura actual presenta brechas significativas. La mayoría de los modelos desarrollados no integran datos en tiempo real obtenidos mediante sensores, ni incorporan técnicas de aprendizaje automático para la calibración dinámica de sus parámetros. Además, suelen omitir la medición directa de gases de efecto invernadero no controlados, lo que restringe su aplicabilidad en condiciones operativas reales. Estas limitaciones evidencian la necesidad de enfoques híbridos que combinen modelamiento mecanístico, monitoreo continuo y análisis predictivo, con el fin de desarrollar sistemas más precisos, adaptativos y contextualizados para la estimación y gestión de emisiones en procesos de bioconversión.


\section{Avances en IoT y AAA para la gestión sostenible de residuos}
La incorporación de tecnologías basadas en IoT y AAA ha demostrado tener un impacto tangible en la sostenibilidad de los sistemas de gestión de residuos orgánicos. En particular, \textcite{galan-diazCarbonWaterFootprint2024} documentan cómo la implementación de sensores y automatización en granjas de \textit{BSF} permitió una reducción significativa de la huella de carbono y agua, gracias al control preciso de variables operativas como la temperatura, la humedad y el consumo energético. Este enfoque no solo mejora la eficiencia del proceso de bioconversión, sino que también permite una trazabilidad más rigurosa del sistema. De forma complementaria, \textcite{jenkinsProcessingPoultryManure2025} evidencian que la integración de tecnologías inteligentes en el tratamiento de estiércol avícola contribuye a una disminución considerable en las emisiones de CO$_2$ y N$_2$O en suelos agrícolas, lo cual refuerza el papel de estas herramientas como catalizadores para la mitigación del cambio climático. En conjunto, estos estudios muestran que la automatización inteligente no solo optimiza los procesos productivos, sino que también fortalece la dimensión ambiental de los sistemas de tratamiento de residuos, consolidándolos como parte de una estrategia integral hacia la circularidad y la neutralidad en carbono.

\section{Brechas tecnológicas y oportunidades de investigación}

A pesar de los avances en tecnologías aplicadas a la gestión de residuos orgánicos, la implementación de sistemas inteligentes en procesos de bioconversión con \textit{BSF} aún enfrenta desafíos importantes. Entre las principales brechas se encuentran: (1) la escasa disponibilidad de sensores con IoT precisos, duraderos y calibrables para medir CO$_2$ y CH$_4$ en entornos de alta humedad; (2) la falta de estandarización en protocolos y formatos que limite la interoperabilidad entre dispositivos y plataformas digitales; (3) la limitada capacidad de los modelos de aprendizaje automático para adaptarse a la complejidad y variabilidad de sistemas biológicos; (4) la ausencia de metodologías robustas de Medición, Reporte y Verificación (MRV); y (5) las dificultades para escalar estas soluciones en contextos rurales o con baja infraestructura \parencite{kamilaris2017review,herrinCollateralDataQuality2021,ahmedAgricultureClimateChange2020,vanIntegrationInternetofThingsSustainable2022}.

Estas brechas representan, al mismo tiempo, oportunidades estratégicas para el desarrollo de sistemas integrales que combinen sensores con IoT especializados, algoritmos de aprendizaje automático adaptativos y plataformas digitales escalables. Se plantea una arquitectura compuesta por cuatro módulos: (i) red de sensores para monitoreo ambiental continuo; (ii) plataforma de análisis en la nube para almacenamiento y procesamiento de datos; (iii) modelos de AAA -como redes neuronales- para el ajuste dinámico de variables y predicción de emisiones; y (iv) una interfaz de visualización y control para la toma de decisiones en tiempo real. Esta solución apunta a mejorar la trazabilidad, reducir costos operativos y facilitar la implementación en escenarios diversos, posicionando a la bioconversión con \textit{BSF} como una herramienta tecnológicamente viable para aportar en la mitigación del cambio climático y la economía circular.

\section{Brechas en el modelamiento matemático de la bioconversión con BSF}

Aunque el modelamiento matemático ha sido fundamental para describir procesos como el crecimiento larval y la conversión de sustrato en biomasa \parencite{padmanabhaComprehensiveDynamicGrowth2020,eriksenDynamicModellingFeed2022,prasetyaGrowthKineticsModel2021}, persisten importantes brechas que limitan su aplicación en escenarios reales. La mayoría de los modelos disponibles operan bajo condiciones estáticas de laboratorio y omiten variables ambientales críticas -como temperatura, humedad del sustrato, oxigenación y composición química- que influyen directamente en la generación de gases de efecto invernadero (GEI) como CO$_2$ y CH$_4$ \parencite{bekkerImpactSubstrateMoisture2021a,rossiEstimatingDynamicsGreenhouse2024}. Esta falta de contextualización reduce su capacidad predictiva y restringe su uso a escala piloto o industrial.

Otra brecha clave es la débil integración con datos experimentales en tiempo real. En sistemas dinámicos, la ausencia de retroalimentación desde sensores con IoT limita la capacidad de los modelos para ajustarse a fluctuaciones operativas, afectando su utilidad en condiciones cambiantes como variaciones súbitas de temperatura o humedad. Asimismo, se observa una gran heterogeneidad metodológica en la validación de modelos: desde balances de masa hasta enfoques empíricos, muchos trabajos carecen de estandarización en métricas, unidades o protocolos de evaluación, lo que dificulta la comparación y escalabilidad de los resultados \parencite{dienerConversionOrganicMaterial2009,fuhrmannComprehensiveIndustryrelevantBlack2025}.

Estas limitaciones abren oportunidades de investigación orientadas al desarrollo de modelos dinámicos basados en ecuaciones diferenciales ordinarias (EDO), integrados con sensores con IoT y algoritmos de aprendizaje automático capaces de recalibrarse de forma adaptativa \parencite{kobelskiModelbasedProcessOptimization2024a,liuOptimizingDataPipelines2023}. Una aproximación híbrida e interdisciplinaria permitiría representar con mayor fidelidad las dinámicas del sistema, mejorar la predicción de emisiones y fortalecer la toma de decisiones operativas en contextos de bioconversión con \textit{BSF}.

